\chapter{Das Dateiformat \texorpdfstring{\datei{.rttf}}{.rttf}}\label{kap:dateiformat}
\index{Dateiformat}\index{rttf}

Dieser Anhang beschreibt, was in einer Dokumentdatei steht. Für die tägliche Arbeit
brauchen Sie ihn nicht; er hilft, wenn Sie Dateien mit eigenen Programmen auswerten oder
erzeugen, eine Datei von Hand reparieren oder verstehen wollen, warum eine Datei groß ist.

\section{Überblick}

Dokumente werden mit der Endung \datei{.rttf} (\enquote{Raytatouille Fusion file})
gespeichert. Dateien mit der Endung \datei{.nfcad} aus der Zeit vor Version~1.0 haben
denselben Inhalt und werden weiterhin geöffnet; \menu{File,Save As...} bietet sie als
\datei{.rttf} an. Intern heißt das Format weiterhin \code{"nfcad"}.

\begin{itemize}
  \item \textbf{Kodierung:} UTF-8-JSON, also eine Textdatei, die sich mit jedem Editor
    lesen lässt.
  \item \textbf{Sicheres Schreiben:} Die Datei wird zuerst vollständig unter einem
    temporären Namen geschrieben und dann in einem Schritt ersetzt. Ein Absturz während des
    Speicherns hinterlässt daher nie eine halbe Datei.
  \item \textbf{Einheiten:} Millimeter.
  \item \textbf{Geometrie exakt:} Gespeichert werden Grade, Knotenvektoren, Kontrollpunkte
    und Gewichte der NURBS -- keine Dreiecksnetze. Beim Öffnen entsteht exakt dieselbe
    Geometrie.
  \item \textbf{Version:} Aktuell ist Version~2; Dateien der Version~1 werden weiterhin
    gelesen. Neuere Versionen, als das Programm kennt, werden mit einer Meldung abgelehnt; das
    geöffnete Dokument bleibt dann unverändert.
\end{itemize}

\section{Aufbau}

Eine Datei ist ein JSON-Objekt mit diesen Abschnitten (gekürzt):

\begin{lstlisting}
{
  "format": "nfcad",
  "version": 2,
  "units": "mm",
  "layers":   [ { "id": 0, "name": "Default", "color": "#bec6d2", "visible": true },
                { "id": 1, "name": "Mounts", "color": "#78a0dc", "visible": true,
                  "locked": true } ],
  "current_layer": 1,
  "groups":   [ { "id": 1, "name": "objective", "parent": 0, "visible": true } ],
  "entities": [ { "id": "1", "name": "part_0_0", "layer": 0, "visible": true,
                  "color": "#ff0000",
                  "group": 1,
                  "attributes": { "rtt.role": "mirror", "sim.inactive": true },
                  "geometry": { "type": "shell", "faces": [ ... ] } } ],
  "plugins":  { "raytracer": { "params": { ... } } },
  "workflow": { "model": { "format": "rtt-fusion-graph", "version": 2, ... }, ... },
  "results":  [ { "key": "workflow:Kamerabild", "name": "Kamera", "kind": "image",
                  "stamp": "3f1c...", "image": { ... } } ]
}
\end{lstlisting}

\begin{referenz}{@{}l>{\raggedright\arraybackslash}X@{}}
\toprule
\textbf{Abschnitt} & \textbf{Inhalt} \\
\midrule\endhead
\param{format}, \param{version}, \param{units} & Kennung \code{"nfcad"}, Formatversion,
  Einheit \code{"mm"}. \\
\param{layers} & Die Ebenen mit Nummer, Name, Farbe, Sichtbarkeit und -- optional, nur wenn
  gesetzt -- \param{locked} für eine gesperrte Ebene (\autoref{kap:organisieren}). \\
\param{current\_layer} & Optional: die aktuelle Ebene, auf der neue Objekte entstehen; fehlt
  sie (oder ist sie unbekannt), ist es \gui{Default} (0). Wie \param{locked} wird sie nur
  geschrieben, wenn sie vom Normalfall abweicht; ältere Programmversionen ignorieren beides. \\
\param{groups} & Die Gruppen des Modellbaums; \param{parent} ist die übergeordnete Gruppe,
  0 bedeutet oberste Ebene. \\
\param{entities} & Die Objekte: Nummer (als Zeichenkette, 64 Bit), Name, Ebene,
  Sichtbarkeit, optional eigene Farbe (sonst die der Ebene), Gruppe, freie Attribute und die
  Geometrie. \\
\param{plugins} & Die Werte der Parameterformulare je Plugin, nach Plugin-Kennung.
  Einträge von Plugins, die das Programm nicht kennt, bleiben beim Speichern unverändert
  erhalten. \\
\param{workflow} & Der Node-Workflow als Modell (\code{rtt-fusion-graph} Version~2): Blöcke
  mit Ports, Verbindungen, Äste, Gruppen mit Varianten und eingebettete Dateien (zum Beispiel
  Glaskataloge), dazu Diagnosen und Plan (\autoref{kap:workflow}). Ältere Dokumente enthalten
  statt des Modells einen Graphen aus Sockets (\param{graph}, Version~1); er wird beim Öffnen
  umgewandelt (\autoref{sec:workflow-alte-form}). Version~1 wird nicht mehr geschrieben. \\
\param{results} & Die gespeicherten Ergebnisse der Anzeigefenster (siehe unten). \\
\bottomrule
\end{referenz}

\section{Geometrie}\index{NURBS}

\begin{referenz}{@{}l>{\raggedright\arraybackslash}X@{}}
\toprule
\textbf{\param{type}} & \textbf{Felder} \\
\midrule\endhead
\param{point} & \param{point}: $[x, y, z]$ \\
\param{curve} & \param{degree}, \param{knots}, \param{points} (flache Liste $x, y, z, \dots$),
  \param{weights} (fehlt, wenn alle Gewichte 1 sind) \\
\param{surface} & \param{degree\_u}, \param{degree\_v}, \param{knots\_u}, \param{knots\_v},
  \param{num\_u}, \param{num\_v}, \param{points} (zeilenweise, Index $i$ entlang $u$),
  \param{weights} \\
\param{trimmed\_surface} & \param{surface}, \param{outer} und \param{holes} (Listen von
  2D-Kurven im Parameterraum mit $z = 0$; fehlen bei ungetrimmten Flächen), \param{label} \\
\param{shell} & \param{faces}: Liste getrimmter Flächen. Die Topologie (Kanten, Ecken) wird
  beim Laden durch Vernähen neu aufgebaut. \\
\bottomrule
\end{referenz}

\section{Attribute}\index{Attribut}

Attribute sind freie Schlüssel-Wert-Paare (Zeichenketten, Zahlen, Wahrheitswerte) an einem
Objekt. Plugins verwenden Schlüssel mit Namensraum, etwa \param{rtt.role} für die Rolle einer
Fläche im Raytracer. Das Attribut \param{sim.inactive} mit dem Wert \code{true} nimmt ein
Objekt von allen Berechnungen aus, ohne es zu löschen oder auszublenden. Attribute lassen sich
im \gui{Properties}-Dock bearbeiten (\autoref{kap:organisieren}).

\section{Gespeicherte Ergebnisse}\index{Ergebnisse!gespeichert}\index{Fingerabdruck}

Der Abschnitt \param{results} enthält das jeweils neueste Ergebnis jedes Anzeigefensters mit
Schlüssel, zum Beispiel \code{workflow:<Block>} für die Anzeigeblöcke des Node-Workflows:
Bilder, Kurven, Tabellen und Balken. Zu jedem Ergebnis gehört ein \emph{Fingerabdruck}
(\param{stamp}) dessen, woraus es berechnet wurde.

\begin{itemize}
  \item Zahlenfelder stehen als Base64 ihrer Little-Endian-Bytes in der Datei: 32-Bit-Gleitkomma
    für Bildwerte, 64-Bit-Gleitkomma für Kurven und Balken, Bytes für Farben.
  \item Beim Öffnen erscheinen sie als Ergebnissatz \gui{Saved results}.
  \item Der Workflow vergleicht die Fingerabdrücke mit dem, was er jetzt berechnen würde. Ein
    unveränderter Anzeigeblock gilt als aktuell (grüner Haken), und \gui{Run workflow} rechnet
    ihn nicht noch einmal (\enquote{(unchanged)} in der Statuszeile).
  \item Zwischenergebnisse (Strahlenbündel, Felder) werden nicht gespeichert -- sie wären
    groß und lassen sich jederzeit neu berechnen.
  \item Programme, die den Abschnitt nicht kennen, überspringen ihn.
\end{itemize}

\begin{tipp}
Wenn eine Datei unerwartet groß ist, liegt das meist an gespeicherten Bildern hoher Auflösung
oder an eingebetteten Glaskatalogen. Beides ist gewollt: Die Datei lässt sich so auf einem
anderen Rechner ohne die Kataloge öffnen und zeigt die Ergebnisse sofort.
\end{tipp}

\section{Geplante Erweiterungen}

Für spätere Versionen vorgesehen sind binäre Blöcke für große Kontrollnetze, gespeicherte
B-Rep-Topologie statt erneuten Vernähens und andere Einheiten als Millimeter.
